% Shared preamble for labbook.tex and worklog.tex.
%
% TEMPLATE FILE — replaced on update. Do not edit; if something is missing here,
% it is a FEEDBACK.md entry.
%
% Deliberately minimal. Figures are produced by the researcher's own plotting
% code and included as image files; no plotting package is loaded and no house
% style is imposed on that code.

\usepackage[T1]{fontenc}
\usepackage[utf8]{inputenc}
\usepackage{amsmath,amssymb,mathtools}
\usepackage{graphicx}
\usepackage{booktabs}
\usepackage{xcolor}
\usepackage{geometry}
\usepackage[hidelinks]{hyperref}

\geometry{margin=1in}
\hypersetup{colorlinks=true,linkcolor=blue,citecolor=blue,urlcolor=blue}
\graphicspath{{figures/}}

% Long run ids and file paths are the main source of overfull lines; let TeX
% break typewriter text rather than run into the margin.
\usepackage[htt]{hyphenat}
\emergencystretch=3em
\tolerance=1500

% --------------------------------------------------------------------------- %
% Macros used by entries. See labbook/00_rules.tex for what each one means.
% --------------------------------------------------------------------------- %

% \entryhead{date}{title}{what this was about, one sentence}
\newcommand{\entryhead}[3]{%
  \section{#1 --- #2}%
  \noindent\textit{#3}\par\medskip%
}

% Provenance of a piece of code or a figure.
\newcommand{\fromcode}[1]{\texttt{#1}}
\newcommand{\atcommit}[1]{\texttt{#1}}

% A transcription of handwritten work, reproduced as written.
\newenvironment{transcription}%
  {\par\medskip\noindent\textbf{Transcribed as written:}\par\nopagebreak\medskip}%
  {\par\medskip}

% Something in the handwriting that could not be read.
\newcommand{\unreadable}[1]{\textcolor{red}{\texttt{[?~#1]}}}

% A flagged possible error, and what was decided about it.
% \flag{date}{what looks wrong}{evidence}{proposal}{decision}
\newcommand{\flag}[5]{%
  \par\medskip\noindent\fbox{\parbox{0.95\linewidth}{%
    \textbf{Flagged #1.} #2\par\smallskip
    \textit{Evidence:} #3\par\smallskip
    \textit{Proposed:} #4\par\smallskip
    \textbf{Decision:} #5}}\par\medskip%
}

% Status of a statement, where it matters.
\newcommand{\unchecked}{\textcolor{red}{[unchecked]}}
\newcommand{\checked}[1]{\textcolor{teal}{[checked #1]}}
